% !PRINTED 5
\weberpage{5}

In the permutation groups considered in the first volume one will easily recognise the marks of the general concept of a group.

We now draw from this definition first some quite general consequences.

By 4. there is for every given $b$ \gesperrt{an} element $e$ in $P$ that satisfies the condition
\weq{1}{eb = b}
and this $e$ is independent of $b$; for from (1) it follows, for every $c$,
\[
ebc = bc,
\]
and $bc$ can by 4. denote every element in $P$. Likewise there is an element $e'$ that for every $b$ satisfies the condition
\weq{2}{be' = b}
This element $e'$, however, is not different from $e$; for if we set $b = e'$ in (1) and $b = e$ in (2), it follows that
\[
e e' = e', \, . \, e e' = e, % [sic: the source prints a full stop between the two identities]
\]
hence
\[
e = e'.
\]

The element $e$ changes nothing when it is compounded with any elements of $P$, and is called the \gesperrt{unit of the group}.

In many cases it can, without misunderstanding, be denoted simply by ``1''.

To every element $a$ there is by 4. a definite element $a^{-1}$ that satisfies the condition
\weq{3}{a^{-1} a = e}
From (3), (1) and (2) it follows that
\[
a^{-1} a a^{-1} = e a^{-1} = a^{-1} = a^{-1} e,
\]
and consequently, by 3.,
\weq{4}{a a^{-1} = e.}

The two elements $a$, $a^{-1}$ are said to be \gesperrt{opposite} or \gesperrt{reciprocal} to one another.

In special cases the commutative law can also hold in the composition of the elements of a group $P$; that is, for any two elements $a$, $b$ of the group
\[
a b = b a
\]
can hold.
